\documentclass[a4paper,12pt]{article}
\usepackage[english,russian]{babel}
\usepackage{fontspec}
\defaultfontfeatures{Ligatures={TeX},Renderer=Basic}
\setmainfont[Ligatures={TeX,Historic}]{Calibri Light}
\setsansfont{Calibri Light}
\setmonofont{Consolas}

\usepackage{indentfirst}
\frenchspacing

% Для математики
\usepackage{amssymb,amsmath}
\setcounter{MaxMatrixCols}{20}
\parindent=24pt
\parskip=0pt
\tolerance=2000

% Для настройки размера страницы
\usepackage{geometry}
\geometry{
	a4paper,
	total={170mm,257mm},
	left=20mm,
	top=20mm,
}

% Для вставки графики
\usepackage{graphicx}

% Для таблиц
\usepackage{longtable}
\usepackage{multirow}
\usepackage{tabu}
\usepackage{colortbl}

% Создаем команду, чтобы переносить текст на новую строку внутри таблицы
\newcommand{\tcell}[2][c]{\begin{tabular}[#1]{@{}c@{}}#2\end{tabular}}

% Пакет для списков
\usepackage[ampersand]{easylist}

% Для контура вокруг текста
\usepackage[outline]{contour}

% Пакеты от tikz
\usepackage{tikz}
\usepackage{graphics}
\usepackage{pgfplots}
\usepackage{pgfplotstable}
\usepackage{xcolor}
\usetikzlibrary{calc}
\usetikzlibrary{through}
\usetikzlibrary{intersections}
\usetikzlibrary{patterns}
\usetikzlibrary{scopes}
\usetikzlibrary{decorations.pathreplacing}
\usetikzlibrary{arrows.meta}

% Для цветных таблиц
\usepackage{tcolorbox}
\tcbuselibrary{skins}
\tcbuselibrary{breakable}
\tcbuselibrary{minted}
\usemintedstyle{vs}

% Для подсветки кода
\usepackage{minted}

% Задаем цвет номеров строк подсветки кода
\renewcommand{\theFancyVerbLine}{\sffamily\textcolor[rgb]{1, 1, 1}{\scriptsize\arabic{FancyVerbLine}}}

% Определяем новую команду для красивой вставки кода в рамочке и с прочими прелестями
\newcommand{\mycodeinput}[3]{
\begin{tcolorbox}[
	colback=black!5!white,
	colframe=black!30!white,
	boxrule=0.5pt, 
	listing only,
	left=-0.5mm,
	leftrule=6mm,
	arc=2mm, outer arc=2mm,
	top=2pt,
	bottom=0pt,
	enhanced,
	breakable,
	title={#3},
	coltitle=black, 
	fonttitle=\bfseries\ttfamily\large,
	attach boxed title to top center={yshift=-3mm,yshifttext=-1mm},
	boxed title style={
		enhanced,
		nobeforeafter,
		tcbox raise base,
		boxrule=0.4pt,
		top=0.5mm,
		bottom=0.5mm,
		right=0mm,
		left=4mm,
		arc=1pt,
		boxsep=3pt,
		before upper={\vphantom{dlg}},
		colframe=black!30!white,
		colback=black!10!white,
		overlay={
			\begin{tcbclipinterior}
			\fill[black!30!white]
				(frame.south west)
					rectangle node[text=white,font=\sffamily\bfseries\footnotesize,rotate=90] {FILE} 
				([xshift=4mm]frame.north west);
			\end{tcbclipinterior}
		}
	},
]
\inputminted[
	breaklines,
	breakanywhere=true,
	autogobble,
	linenos,
	numbersep=1mm,
	mathescape, 
	fontsize=\footnotesize, 
	tabsize=4
]{#1}{#2}
\end{tcolorbox}
}

% Определяем новую команду для вставки кода прямо в тексте
\newtcbox{\inlinecodetable}{on line,
arc=2pt,colback=gray!10!white,colframe=gray!50!black,
before upper={\rule[-3pt]{0pt}{10pt}},boxrule=0.5pt,
boxsep=0pt,left=2pt,right=2pt,top=2pt,bottom=0pt}
\newcommand{\mycodeinline}[2]{\inlinecodetable{\mintinline{#1}{#2}}}

% Определяем новую команду для создания титульного листа
\newcommand{\mytitlepage}[9]{
\begin{center}
\hfill \break
\Large{Министерство науки и высшего образования Российской Федерации}\\
\hfill \break
\large{Федеральное государственное бюджетное образовательное учреждение высшего образования}\\ 
\normalsize{\textbf{«НОВОСИБИРСКИЙ ГОСУДАРСТВЕННЫЙ ТЕХНИЧЕСКИЙ УНИВЕРСИТЕТ»}}\\
\hfill \break
\includegraphics{nstu_logo.eps} \\
\hfill \break
\large{Кафедра #1}\\
\hfill \break
\large{Лабораторная работа №#2\\по дисциплине <<#3>>}\\
\hfill \break
\hfill \break
\Large{\textbf{#4}}\\
\hfill \break
\hfill \break 
\normalsize{\begin{tabular}{cllp{1.5cm}p{1.5cm}}
\multirow{5}{*}[0.75cm]{\includegraphics[scale=0.5]{fami_logo.eps}}
& \textbf{Факультет:} & ПМИ & & \\[1.25ex]
& \textbf{Группа:} & #5 & & \\[1.25ex]
& \textbf{Студент:} & #6 & & \\[1.25ex]
& \textbf{Вариант:} & #7 & & \\[1.25ex]
& \textbf{Преподаватель:} & #8 & & \\[1.25ex]
\end{tabular}} \\
\hfill \break
\hfill \break
\hfill \break
\hfill \break
\hfill \break
\large{Новосибирск\\#9}
\end{center}
\thispagestyle{empty}
\newpage 
\setcounter{page}{1}
}

\begin{document}

\mytitlepage{Прикладной математики}{3}{Комьютерная графика}{Трассировка лучей}{ПМ-63}{Шепрут И.И.}{-}{Задорожный А.Г.}{2019}

\section{Цель работы}

Ознакомиться с основными аспектами метода трассировки лучей.

\section{Постановка задачи}

\noindent\normalsize{\begin{easylist}
\ListProperties(Hang1=true)
& Считывать из файла (в зависимости от варианта):
&& Тип объекта.
&& Координаты и размер объектов.
&& Параметры материала объектов.
& Выполнить трассировку первичных лучей.
& Добавить зеркальную плоскость и учесть отраженные лучи.
& Предусмотреть возможность включения/исключения объектов.
& Предусмотреть возможонсть изменения положения источника света.
\end{easylist}}

\section{Реализованные функции}

\noindent\normalsize{\begin{easylist}
\ListProperties(Hide=100, Hang1=true, 
	Progressive*=.5cm,%	
	Style1*=\textbullet ,%
	Style2*=\textopenbullet )
& \textbf{Два типа рендеринга.}
&& Ray Tracing --- обычная трассировка лучей, которая может отобразить картинку с малым количеством семплов. Трейсинг лучей останавливается при достижении материала с цветом, и рендерит вторичные лучи, только если попадает на преломляющую или зеркальную поверхность.
&& Path Tracing ---  методика рендеринга в компьютерной графике, которая стремится симулировать физическое поведения света настолько близко к реальному, насколько это возможно. Трейсинг лучей останавливается только при достижении предела отражений, или при достижении источника света.
& \textbf{Полиморфизм.}
&& Полиморфизм объектов --- имеется два типа объектов: object и shape.
&&& Object --- объект, котороый неразрывно связан со своим собственным материалом, например: небо, cubemap, текстурированный полигон.
&&& Shape --- объект, которому можно задать любой имеющийся материал, например, это объекты: сфера, цилиндр, полигон, треугольник, портал, контур (состоит из сфер и цилиндров).
&& Полиморфизм камер --- можно рендерить картинку совершенно различными камерами, которые наследуются от абстрактного класса базовой камеры и реализуют виртуальные методы. Таким образом уже реализованы камеры: перспективная, ортогональная, 360. Так же можно реализовать другие камеры, например: проекцию Панини, камера для создания cubemap.
&& Полиморфизм материалов --- любому объекту можно задать любой материал на его поверхности, например реализованы такие материалы, как: рассеивающий, отражающий, стекло, освещение. Аналогично для реализации материала надо наследоваться от базового класса материала и реализовать метод обработки луча материалом.
& \textbf{Точечные источники света} поддерживаются на уровне рендерера.
& \textbf{Рендеринг порталов.} Портал является shape, потому что можно задать какой материал будет на двух его задних частях. Он сам задается как полигон и две системы координат.
&& Поддерживается телепортирование освещения от точечных источников света через порталы. При этой телепортации учитывается, что некоторая часть света может не пройти через портал.
& \textbf{Считывание произвольной сцены из многоугольников, с анимацией, из \texttt{json} файла.}
& \textbf{Поддержка текстур.} Реализуется эта поддержка в классе текстурированного полигона.
& Рендеринг произвольных полигонов.
& Для path tracing камерой поддерживается симуляция прохождения луча через диафргаму, поэтому можно получить эффект depth of field: когда камера на чем-то сфокусирована, а остальная часть мира размыта.
& Рендеринг происходит в многопоточном режиме.
& Имеется возможность одновременно с обычным изображением рендерить изображение глубины.
& Имеется возможность сохранять отрендеренные изображения в png.
& При рендеринге пиксели выбираются случайным образом, чтобы максимально точно предсказывать итоговое время рендеринга, которое выводится на экран.
\end{easylist}}

\section{Пример реализованных функций}

\newcommand{\img}[2]{
\begin{center}
\includegraphics[width=#2\textwidth]{#1}	
\end{center}
}

\newcommand{\imgtwo}[4]{
\noindent\begin{center}
\includegraphics[width=#2\textwidth]{#1}
\includegraphics[width=#4\textwidth]{#3}
\end{center}
}

\subsection{Различие Ray tracing от Path tracing}

\imgtwo{img/standard_scene_ray_perspective.png}{0.49}{img/standard_scene_path_perspective.png}{0.49}

\imgtwo{img/standard_scene_2_ray_perspective.png}{0.49}{img/standard_scene_2_path_perspective.png}{0.49}

\subsection{Различные виды камер}

\subsubsection{Ортогональная}

\img{img/standard_scene_ray_orthogonal.png}{0.7}

\img{img/standard_scene_2_ray_orthogonal.png}{0.7}

\subsubsection{Камера на 360 градусов}

Эта камера создает такое изображение, которое запечатляет всю обстановку вокруг. Далее, по этому изображению можно получить сколько угодно изображений перспективной проекции, которые смотрят с разных углов при помощи различных программ для просмотра такого рода изображений.

Изображения в 360 градусов были просмотрены с помощью \href{https://www.chiefarchitect.com/products/360-panorama-viewer/}{www.chiefarchitect.com}.

\img{img/standard_scene_ray_360.png}{1}

\img{img/screen_360_1.png}{1}

\img{img/standard_scene_2_ray_360.png}{1}

\img{img/screen_360_2.png}{1}

\subsection{Телепортация освещения от точечного источника света через портал}

\img{img/portal_light_test1_1.png}{1}

\imgtwo{img/portal_light_test1_2.png}{0.49}{img/portal_light_test1_3.png}{0.49}

\subsection{Изображение глубины каждого пикселя}

\img{img/standard_scene_ray_perspective_depth.png}{0.7}

\subsection{Рендеринг произвольных полигонов}

\img{img/polygon_test.png}{0.7}

\subsection{Рендеринг сцен из предыдущей работы, заданных в json}

\imgtwo{img/oblique_portal.png}{0.49}{img/volumetric_portal_with_textures.png}{0.49}

\imgtwo{img/scale_portal.png}{0.49}{img/volumetric_portal_with_textures_depth.png}{0.49}

\subsection{Пример предсказания времени рендеринга}

\img{img/log.png}{0.7}

\section{Код программы}

\subsection{Файлы заголовков}

% \mycodeinput{c++}{code/fragment.h}{fragment.h}
% \mycodeinput{c++}{code/framebuffer.h}{framebuffer.h}
% \mycodeinput{c++}{code/opengl_common.h}{opengl\_common.h}
% \mycodeinput{c++}{code/plane.h}{plane.h}
% \mycodeinput{c++}{code/scene_reader.h}{scene\_reader.h}
% \mycodeinput{c++}{code/shader.h}{shader.h}

\subsection{Исходные файлы}

% \mycodeinput{c++}{code/main.cpp}{main.cpp}
% \mycodeinput{c++}{code/fragment.cpp}{fragment.cpp}
% \mycodeinput{c++}{code/framebuffer.cpp}{framebuffer.cpp}
% \mycodeinput{c++}{code/opengl_common.cpp}{opengl\_common.cpp}
% \mycodeinput{c++}{code/plane.cpp}{plane.cpp}
% \mycodeinput{c++}{code/scene_reader.cpp}{scene\_reader.cpp}
% \mycodeinput{c++}{code/shader.cpp}{shader.cpp}

\subsection{Шейдеры}

% \mycodeinput{c++}{code/draw.fragment.glsl}{draw.fragment.glsl}
% \mycodeinput{c++}{code/draw.vertex.glsl}{draw.vertex.glsl}
% \mycodeinput{c++}{code/drawpoly.fragment.glsl}{drawpoly.fragment.glsl}
% \mycodeinput{c++}{code/drawpoly.vertex.glsl}{drawpoly.vertex.glsl}
% \mycodeinput{c++}{code/merge.fragment.glsl}{merge.fragment.glsl}
% \mycodeinput{c++}{code/merge.vertex.glsl}{merge.vertex.glsl}

\end{document}